\section{Related Work}

\textbf{Mixing Methods.} The offline mixing schema, also described as ``function fitting-based offline methods'' in~\citet{liu2025rethinkingdatamixturelarge} has been extensively used in existing mixing methods. In Table~\ref{tab:mixing_design_choices}, we provide an overview of several offline mixing methods and how they address the key design choices we study. Some methods use explicit parametric regression models~\citep{que2024dcptlawdomainspecificcontinual}, while others use nonparametric approaches like LightGBM and Gaussian Processes~\citep{liu2025regmixdatamixtureregression, chen2025admirebayesoptaccelerateddatamixture}. Some methods explicitly model the role of proxy model size or training budget~\citep{ge2025bimixbivariatedatamixing, ye2025datamixinglawsoptimizing, kang2025autoscalescaleawaredatamixing} while others assume direct generalization from proxy models to target models~\citep{held2025optimizingpretrainingdatamixtures}. There are also several methods that build on top of RegMix~\citep{liu2025regmixdatamixtureregression}, exploring how it can be augmented with better domains~\citep{wettig2025organizewebconstructingdomains}, iterative swarms~\citep{diao2025climbclusteringbasediterativedata}, or more features~\citep{belenki2025optimizingpretrainingdatamixtures}. Our work examines key design choices for the offline schema that underlies all these methods. 


Aside from the offline schema, DoReMi~\citep{xie2023doremioptimizingdatamixtures} and DoGE~\citep{fan2024dogedomainreweightinggeneralization} determine a mixture by using one or two proxy runs that dynamically explore the mixture weight space; this is in contrast to using many proxy runs with static mixes. While these approaches have proven effective in some settings, previous analysis~\citep{chen2025aioliunifiedoptimizationframework} has suggested some suboptimality in how the dynamic update rules are constructed, and the offline schema is generally considered simpler to implement.


Online mixing methods~\citep{chen2023skillitdatadrivenskillsframework, jiang2024adaptivedataoptimizationdynamic, albalak2023efficientonlinedatamixing} adjust the mixture weights throughout the final training run. Rather than exploring the mixture space and learning from it offline, these methods do this on the fly during training. Note that our notion of dynamics, which is over the domain set and during LM development (i.e., before the final training run), is different from the dynamic aspect of online mixing methods.

Existing mixture literature largely assumes fixed domain sets. The recent work Chameleon~\citep{xie2025chameleonflexibledatamixingframework} addresses adaptability to domain changes by computing domain weights from learned embeddings using kernel ridge leverage scores, which allows direct transfer to new data without proxy retraining. Chameleon focuses on adding new domains and computes weights directly from domain embeddings without the swarm-based regression approach of offline methods. In contrast, our mixture reuse approach explicitly handles various domain update operators (add, remove, partition, and revise) and is designed to work with any method that follows the offline mixing schema. Moreover, our approach provides a theoretical and empirical analysis of when and why existing mixes can be reused.



\textbf{Data-constrained settings.} Several works have studied how to train language models under data constraints. UniMax~\citep{chung2023unimaxfairereffectivelanguage} proposes an allocation algorithm for multilingual pretraining that distributes a fixed token budget to maximize uniform coverage across languages while capping repetitions to avoid overfitting on low-resource languages. \citet{muennighoff2025scalingdataconstrainedlanguagemodels} find that up to 4 epochs of repetition yields negligible degradation and propose modified scaling laws that account for diminishing returns of repeated tokens. Our work explicitly integrates repetition constraints directly into data mixing, enabling principled allocation of limited data budgets while producing a mix that yields strong downstream performance.

\textbf{LM Data Development.} Real-world LM development involves iterative refinement of training data. Works like DCLM~\citep{dclm}, Dolma~\citep{soldaini2024dolma}, and FineWeb~\citep{penedo2024fineweb} extensively document the curation processes involved in creating high-quality pretraining corpora, including quality filtering, deduplication strategies, and careful domain selection. Continuous projects like OLMo 1-3~\citep{Groeneveld2024OLMoAT, olmo20242olmo2furious, olmo2025olmo3} and SmolLM 1-3~\citep{allal2024smollm, allal2025smollm2smolgoesbig, bakouch2025smollm3} showcase how training data evolves over time. Motivated by these works, \olmix views data mixing as an ongoing process throughout LM development.